\section{Discussion}
\label{sec:discussion}

\subsection{What the results do and do not show}

We began with the expectation that the architectures would order themselves by
how much of the expert parameterisation is shared: independent experts at one
end, product-key composition in the middle, multilinear factorisations at the
other, with generalisation improving monotonically along that axis. The two
endpoints behave as expected. The middle does not survive contact with the
data, because Tensor-Train sits firmly at the shared end and scales like the
product-key layers rather than like Tucker.

What remains is a weaker and more useful statement. Every architecture we
tested has a window of expert counts in which it generalises, and the
factorisation determines where that window sits and how wide it is.
Independent experts have the narrowest: they work at $N = 16$ and $N = 64$ and
fail completely at $N = 256$. CP has a window with a visible lower edge --- at
$N = 16$ two seeds in three fail --- and no upper edge within our range.
Tucker's window covers the entire range, with its best performance at the top.
MONET-HD degrades steadily and begins to fail at $N = 4096$. Whether adding
experts helps is therefore a question about where a given architecture
currently sits in its own window, not a question about mixtures of experts.

Two mechanisms plausibly set the two edges, and the ablations let us say
something about each rather than speculate.

The upper edge is at least partly the router. Under one-hot inputs the router
keys form a table of $8nP$ parameters addressed directly by the operand, and it
grows linearly in $n$. Section~\ref{sec:res-ablations} shows this is a real
memorisation channel and not a hypothetical one: at $N = 4096$, exempting those
keys from weight decay costs MONET-HD generalisation outright, and at a $50\%$
training fraction the same exemption is the difference between $0.006$ and
$0.977$ test accuracy. The channel is harmless while memorisation is
unattractive and becomes decisive when it is not, which is exactly the regime
large expert counts create.

The lower edge is partly rank starvation in the expert modes, but only on one
side. Our Tucker layer uses an expert-mode rank of $\min(32, n)$, so at $n = 4$
the routing coefficients are compressed into a four-dimensional space and the
experts are nearly indistinguishable; CP at $n = 4$ projects a rank-$128$
factor onto four columns. Both architectures are at their worst there, and both
improve as $n$ rises to meet their rank. Section~\ref{sec:res-rank} tests this
directly by freezing the Tucker rank at $4$, and the improvement disappears:
the layer stops fitting at all at $N = 256$ and then slows from $1800$ to
$2400$ steps as experts are added, where the unfrozen layer stays near $600$.

The same account predicted that Tensor-Train, whose expert-mode bond we held at
rank $16$ throughout the sweep, would be rescued by letting that rank grow. It
is not. Raising the rank leaves the upward trend in expert count intact at
every value tested and makes the layer worse at every expert count, opening a
failure region that widens with rank until, at rank $64$ and $N = 256$, a
layer of over a million parameters memorises the training set and reaches
$0.022$ test accuracy. We state the prediction and its failure because the
failure is the more informative half. The sign of the rank effect turns out to
be a property of the particular factorisation, exactly as the sign of the
expert-count effect is, and our search for a single quantity that sets the
direction for every architecture has now come up empty twice. The upper-edge
account stands as it is, since the router ablation supports it directly.

\subsection{Consequences for routing-score interpretability}

The router result is the one that does not depend on which architecture is
better. Across all six mixture layers, and regardless of whether that layer
generalises sooner or later with scale, the spectral purity of the router keys falls as experts are added, reaching the white-noise floor at the largest expert count each is run at. In
the same layers where the shared basis stays clean --- Tucker at purity
$0.951$, Tensor-Train at $0.976$ on the seeds that behave --- the router has
fallen to $0.048$ and $0.059$. The computation and the routing end up in
different places, and the gap widens with scale.

The obvious reading of that --- routing scores stop being informative at scale
--- is wrong, and Section~\ref{sec:res-mixed} shows why. On the mixed task the
same routers at the same expert count support MONET's criterion perfectly:
factors assigned to a subtask by routing-score skew are exactly the factors
whose removal destroys that subtask and spares the other, with selectivity
$0.99$ at $N = 4096$. A router that scores at the white-noise floor on
frequency content is still a flawless detector of which subtask it is looking
at.

The correct statement is therefore about what kind of structure survives.
Routing scores recover partitions that the input already carries, and miss the
structure that performs the computation; scaling drains the second and leaves
the first. This is a more uncomfortable position for routing-based
interpretability than plain uninformativeness would be. A method that
identifies experts by routing-score skew and validates the identification by
ablation is checking one reading of the partition against another, so the two
agree by construction, and the agreement certifies nothing about mechanism.
Nothing inside such a method signals which of the two situations it is in.
MONET's domains --- Python, chemistry, toxic text --- are partitions of the
kind that is largely readable from the tokens, which is precisely the regime
where the criterion will look most convincing and mean least.

We would not extrapolate to language models on the strength of one algorithmic
testbed. What the testbed establishes is that the dissociation can occur, that
the frequency half of it strengthens with scale, and that the standard
validation procedure cannot detect it. Testing for it in a language model does
not require Fourier analysis: it requires a task where the mechanism is known
independently of the domain label, so that routing-score attribution can be
scored against something other than the partition it was derived from.

\subsection{Limitations}
\label{sec:limitations}

\paragraph{Structural alignment.} The multilinear layers compute a product of a
factor depending on $a$ and a factor depending on $b$, and the Fourier solution
of modular addition has exactly that form. The architecture is aligned with the
target function in a way specific to this task. We measured the size of that
advantage rather than leaving it as a caveat: permuting the router's input so
that its halves are no longer the two operands slows Tucker by $8.5\times$
under interleaving and $2.8\times$ under a random permutation, and makes
MONET-HD fail entirely. Alignment therefore accounts for a substantial part of
what we report, and the correct statement of the result is that a decomposition
whose structure matches the structure of the task both generalises sooner and
yields a sparser representation in the task's own basis. Establishing anything
more general requires tasks whose solutions are not multiplicative in the
operands.

\paragraph{One task, one modulus.} All results are for $c = (a+b) \bmod 113$
with one-hot inputs and a single hidden layer. Nothing here demonstrates
transfer to language models, which is where MONET operates and where the
monosemanticity claims we examine were made.

\paragraph{Mechanisms are partly tested.} The router-table account of the upper
edge is supported by the ablations, and the rank account of the lower edge has
now been tested on both sides: confirmed for Tucker, refuted for Tensor-Train,
which leaves the Tensor-Train trend without a mechanism.

\paragraph{Ranks and widths are not matched.} CP is given rank $128$, Tucker
$32$, Tensor-Train $16$, and every layer an expert dimension of $8$. These were
chosen to put the architectures in a comparable parameter range rather than by
any principle, and Section~\ref{sec:res-comp} shows that at least one of them
is an order of magnitude from binding. Differences we attribute to the
structure of a decomposition may therefore be partly differences in how
generously each one happens to be provisioned.

\paragraph{One training regime, selected on one architecture.} The protocol of
Section~\ref{sec:protocol} was fixed by a sweep over the dense MLP and then
applied unchanged to every Mixture-of-Experts layer, although
Section~\ref{sec:protocol} itself observes that the grokking band depends on
the architecture. Learning rate, weight decay, the Adam betas and the
auxiliary weight are therefore tuned for one member of the comparison and
inherited by the other six. This is the largest confound in the study. It
bears most heavily on the configurations that fail: MONET-VD fails at four of
five expert counts and generalises at one, and we cannot distinguish an
architectural property from a regime that suits it badly at the other four.

\paragraph{Weight decay is not the same regulariser in every format.} An $L_2$
penalty applied to CP factors induces a penalty on the reconstructed tensor
close to a nuclear norm; applied to a Tucker core and its factors it induces a
different one; applied to a chain of Tensor-Train cores, a third. Holding the
coefficient at $1.0$ across architectures therefore does not hold the
regularisation fixed --- it varies it in a way determined by the very
factorisation we are trying to isolate. This is a plausible partial account of
why the sign of the expert-count effect differs between formats, and it is not
one our experiments can separate from the structural account we give.

\paragraph{Seeds are uneven across experiments.} Section~\ref{sec:method}
states that three seeds are used for every configuration carrying a claim, and
that holds for the main grid and for the router-key decay ablation. It does
not hold for the remaining ablations, the rank sweeps of
Sections~\ref{sec:res-rank} and~\ref{sec:res-comp}, or the mixed-task control,
all of which are seed $0$. For architectures whose seed spread reaches $6861$
steps, single-seed differences below roughly $2000$ steps should be read as
suggestive.

\paragraph{Degradation factors are lower bounds.} Failed runs are excluded
from the averages rather than assigned the step budget, so a statement such as
``MONET-HD degrades by $7.4\times$'' compares the seeds that survived at each
expert count. Where some seeds fail at the larger count, as at $N = 4096$, the
true degradation is worse than the figure given.

\paragraph{Frequency sparsity is a correlate, not a law.} Within the mixture
family, narrower frequency budgets accompany earlier generalisation, but the
dense MLP violates the trend, using $38.1$ effective frequencies and still
generalising faster than every product-key variant. CP broadens its basis
steadily with scale while its grokking step falls. We therefore describe the
relationship as a regularity within the family rather than as a mechanism.

\paragraph{Unexplained cases.} MONET-VD fails at $N = 16$, $64$, $1024$ and
$4096$ and generalises only at $N = 256$; we report this rather than omit it,
and have no account of it. One Tensor-Train seed at $N = 4096$ reaches perfect test accuracy with a basis
indistinguishable from noise; Section~\ref{sec:res-rank} shows this is
systematic at high rank and is most likely a limitation of how we extract the
basis for this format rather than a property of the model.

\paragraph{Fixed data partition and single-seed ablations.} Seeds vary
initialisation only; the train/test split is held fixed, which isolates
initialisation variance but says nothing about sensitivity to which pairs are
withheld. Most of the ablations use one seed, so the differences reported there are not
protected against the seed variance that Section~\ref{sec:res-stability}
documents for the product-key layers in particular; the router-key decay ablation, on which the upper-edge account rests, is the exception and is run on three.

The mixed-task control of Section~\ref{sec:res-mixed} is likewise seed $0$
only, and ablates router factors rather than composed experts, which
product-key composition does not allow to be removed individually.

\subsection{Future work}
\label{sec:future}

The most direct continuation is the decomposition this study did not train.
Tucker performs best here, but its core grows as $\prod_n R_n$, which is why
its parameter count still rises steeply with expert count in
Table~\ref{tab:params} while CP's and Tensor-Train's do not.
Section~\ref{sec:res-comp} shows that most of that growth is unused, which
makes the hierarchical Tucker format \citep{grasedyck2010hierarchical} a way of
removing it by construction rather than by a rank constraint imposed after the
fact: instead of one dense core, the modes are arranged on a binary tree with a
small transfer tensor at each internal node, giving parameter count linear
rather than exponential in the order while retaining the pairwise mode
interactions that CP lacks.

Applied to an expert layer the construction is natural. Factor the expert index
into $L$ binary modes, $N = 2^L$, and place them at the leaves of a balanced
tree. Routing then becomes hierarchical: a token descends the tree, making one
binary decision per level, and selects an expert in $O(\log N)$ rather than
$O(\sqrt{N})$ decisions. Two properties follow that bear on the questions
raised here. The internal nodes carry representations at intermediate levels of
granularity, so the hierarchy itself becomes an object of interpretation ---
coarse concepts near the root, fine ones near the leaves --- rather than only
the leaves, as in a flat expert set. And the transfer tensors are shared across
all experts beneath a node, so sharing is graded by depth rather than fixed by
a single global choice. Given that our results turn on where each
architecture's usable window of expert counts sits, an architecture with a
tunable sharing gradient is the natural instrument for asking what sets the
edges of that window.

Three further directions follow from the limitations. Identifying what does set the direction for Tensor-Train, now that rank is
ruled out, together with an extraction of the operand representation from the
whole core chain rather than from the input core alone. Tasks whose solutions are not
multiplicative in the operands, which would separate the effect of
factorisation from the effect of structural alignment. And the same
router-versus-mechanism comparison on a small language model, where the shared
basis cannot be Fourier-analysed but routing-score attribution and ablation can
still be set against each other.

\section{Conclusion}

We asked whether scaling the number of experts in a Mixture-of-Experts layer
makes it more interpretable, on a task where the answer can be checked rather
than illustrated. It does not, and the reason is more specific than we
expected. Adding experts brings generalisation forward for some factorisations
of the expert tensor, pushes it back for others, and destroys it for
independent experts --- with Tensor-Train and Tucker landing on opposite sides
despite sharing their parameters equally thoroughly. Each architecture has a
window of expert counts in which it works, and the factorisation sets where
that window lies.

One result holds across all of them. As experts are added, the router's
frequency content decays to noise in every mixture architecture we tested,
while the shared basis of the best-scaling ones stays clean --- so the
computation and the routing end up in different places, and the gap widens
with scale. The router does not go blank: on a control task whose subtask is
signalled in the input, routing-score attribution still identifies the
responsible factors perfectly at the largest expert count we tested. What it
recovers is the partition the input already carries; what it misses is the
structure that performs the computation, and from inside the method the two
look the same. Whether that dissociation holds in language models is the
question this testbed was built to make askable.